\scanpage{263}
\mspage{75}
% archive sheet number: 263
En voici un exemple typique, en supposant $\mathbb D$ un dérivateur \uline{pointé} (les $\mathbb D(I)$ sont pointées, et les $f^{\unsure{\prime|*}}$ transforment \deleted{\unsure{l'|les}}objets nuls en objets nuls).

Il s'agit de la définition du foncteur \deleted{``\unsure{lacets}''} ``cosuspension'' ou ``\unsure{espaces|espace} de Lacets'' de $\mathcal A$

\[
\Omega:\mathcal A\longrightarrow\mathcal A.
\]

Il est entendu que pour $\deleted{\unsure{N}}M$ de $\mathcal A$, $\Omega\deleted{\unsure{N}}(M)$ doit s'insérer dans un carré commutatif

\diagram{$(C)\qquad$\begin{tikzcd}[ampersand replacement=\&,row sep=large,column sep=large,execute at end picture={
\draw[thin] ([xshift=-1.4em]\tikzcdmatrixname-1-3.west) -- (\tikzcdmatrixname-1-3.north east) -- ([xshift=-1.4em,yshift=.4em]\tikzcdmatrixname-2-3.west) -- (\tikzcdmatrixname-1-3.south east) -- ([xshift=-1.4em]\tikzcdmatrixname-2-3.west) -- (\tikzcdmatrixname-2-3.east);
}]
e\arrow[d]\&\Omega M\arrow[l]\arrow[d]\arrow[r]\&\deleted{e}\arrow[d]\\
M\&e\arrow[l]\arrow[r]\&\deleted{M}
\end{tikzcd}\,,}

où $e$ désigne les éléments unités. Ce carré cependant devrait être substantiel, plus \deleted{\unsure{exactement}}\inserted{précisément}, être un carré substantiel ``\uline{$\mathbb D$-cartésien}''.

Étant entendu qu'un objet de $\mathbb D(I)$ \deleted{\unsure{est|objet}} (carré substantiel de $\mathcal A$) est dit $\mathbb D$-cartésien, s'il est dans l'image essentielle de $\mathbb D(\unsure{\mathbb I_1|I_1})$ par $i_{1*}$, où

\[
i_1:\ \mathbb I_1\lhook\joinrel\longrightarrow I\quad\deleted{\unsure{\mathrm{est}}},
\]

\deleted{\unsure{est défini par l'immersion}} $\mathbb I_1$ étant la \unsure{sous-catég\textsuperscript{ie}|sous-catég\textsuperscript{orie}} ouverte \inserted{de $I$} définie par

\diagram{$\mathbb I_1={}$\begin{tikzcd}[ampersand replacement=\&,row sep=small,column sep=large]
a\&\\
0\arrow[u]\arrow[r]\&b
\end{tikzcd}\,.}
